\chapter{خلاصه اجرایی، واژه‌نامه و منشور اصول}
\label{app:summary_glossary}

این بخش برای دسترسی سریع به مفاهیم کلیدی و اصول اخلاقی طرح تدوین شده است.

\section{مدل مهستان در ۱۰ اصل کلیدی}

\begin{modernbox}{خلاصه‌ی اجرایی}
\begin{enumerate}
    \item \textbf{نهاد محوری:} تشکیل مجلس مهستان با ۳۰۰ نماینده منتخب مستقیم.
    \item \textbf{مشروعیت:} تمام قدرت ناشی از رأی آزاد مردم است.
    \item \textbf{دو قانون اساسی:} تفکیک میان قانون اساسی موقت و نهایی.
    \item \textbf{عدالت انتقالی:} ۵ ستون برای حقیقت، عدالت، غرامت، آشتی و اصلاح.
    \item \textbf{اصلاح نه تخریب:} حفظ ساختار اداری و علمی کشور با تغییر رهبری سیاسی.
    \item \textbf{حقوق بنیادین فوری:} لغو سانسور، حجاب اجباری و آزادی زندانیان از روز اول.
    \item \textbf{تمرکززدایی:} به رسمیت شناختن تنوع قومی و زبانی ایران.
    \item \textbf{اقتصاد شفاف:} حسابرسی بنیادها و بازیابی ثروت‌های غارت‌شده.
    \item \textbf{زمان محدود:} کل دوره‌ی گذار حداکثر ۲ سال به طول می‌انجامد.
    \item \textbf{آمادگی از هم‌اکنون:} دعوت از تمام جریان‌ها برای تدوین پیش‌نویس‌ها.
\end{enumerate}
\end{modernbox}

\section{واژه‌نامه‌ی تخصصی (Glossary)}

\begin{center}
\small
\captionof{table}{واژگان کلیدی گذار دموکراتیک}
\begin{tabular}{R{3.5cm} L{4.5cm} R{6cm}}
\hline
\textbf{اصطلاح فارسی} & \textbf{معادل انگلیسی} & \textbf{توضیح مختصر} \\
\hline
%\endfirsthead
%\hline
%\textbf{اصطلاح فارسی} & \textbf{معادل انگلیسی} & \textbf{توضیح مختصر} \\
%\hline
%\endhead
%\hline
%\endfoot
\textbf{عدالت انتقالی} & \lr{Transitional Justice} & سازوکارهای حقوقی برای رسیدگی به جنایات گذشته \\
\textbf{اصلاح بخش امنیتی} & \lr{Security Sector Reform} & بازسازی نهادهای مسلح تحت نظارت مدنی \\
\textbf{غربال‌گری (\lr{وِتینگ})} & \lr{Vetting} & بررسی سابقه کارکنان برای حذف ناقضان حقوق بشر \\
\textbf{قانون اساسی موقت} & \lr{Interim Constitution} & منشور حقوقی موقت برای مدیریت دوران گذار \\
\textbf{مجلس مؤسسان} & \lr{Constituent Assembly} & مجلس منتخب برای تدوین قانون اساسی دائمی \\
\textbf{نظام نسبی فهرستی} & \lr{List PR system} & توزیع کرسی‌ها بر اساس نسبت آرای هر حزب \\
\textbf{آمبودزمان} & \lr{Ombudsman} & نهاد مستقل رسیدگی به شکایات مردم از دولت \\
\textbf{حقیقت‌یابی} & \lr{Truth-seeking} & فرایند مستندسازی جنایات و شنیدن صدای قربانیان \\
\hline
\end{tabular}
\end{center}

\section{منشور اصول بنیادین غیرقابل‌نقض}

این منشور باید توسط تمامی جریان‌های سیاسی شرکت‌کننده در انتخابات امضا شود.

\begin{principlebox}{حاکمیت مردم}
تنها منبع مشروعیت سیاسی، رأی آزاد مردم است. هیچ فرد، گروه یا ایدئولوژی‌ای حق حاکمیت فراقانونی ندارد.
\end{principlebox}

\begin{principlebox}{جدایی دین از دولت}
دولت در برابر تمام ادیان و عقاید بی‌طرف است. آزادی عبادت و عقیده برای همگان تضمین می‌شود.
\end{principlebox}

\begin{principlebox}{برابری کامل}
تمام شهروندان صرف‌نظر از جنسیت، قومیت، مذهب و گرایش جنسی در برابر قانون برابرند.
\end{principlebox}

\begin{principlebox}{ممنوعیت شکنجه}
شکنجه و مجازات‌های خشن و غیرانسانی تحت هر شرایطی و برای هر کسی مطلقا ممنوع است.
\end{principlebox}

\begin{principlebox}{تمامیت ارضی}
حفظ تمامیت ارضی ایران با پذیرش تکثر فرهنگی و تمرکززدایی اداری-سیاسی.
\end{principlebox}
